\begin{description}
\item[a)]
% FIGURE NEEDED: ray diagram of a telescope objective focusing two ray
% bundles (solid and dashed) separated by angle theta onto points separated
% by S at the focal plane; 2017_Theory_Solution.pdf page 20.
For any two points on the sky separated by small angular distance $\theta$,
at the focal plane of the telescope the distance is $S$
\[ \tan\theta = \frac{S}{f}, \]
\hfill(0.5 marks)\\
where $f$ is the focal length.

For small angle $\theta$
\[ \tan\theta \approx \theta \]
Then plate scale is
\[ \frac{\theta}{S} = \frac{1}{f} \]
\hfill(0.5 marks)
\[ f = \text{F-ratio} \times \text{Diameter} \]
\hfill(1 mark)
\[ \text{Plate scale} = \frac{1}{\text{F-ratio} \times \text{Diameter}} \]
\hfill(1 mark)
\begin{align*}
\text{Plate Scale}
  &= \frac{1}{2.5 \times 40 \times 10^{-2}\,(\mathrm{m})
     \times 10^{3}\,(\mathrm{mm\,m^{-1}})}
     \times \frac{180 \times 3600\,(\mathrm{arcsec})}{\pi}
  && \text{(1 mark)} \\
  &= 206~\mathrm{arcsec/mm} \\
\text{Plate Scale} &= 3.44~\mathrm{arcmin/mm}
  && \text{(1 mark)}
\end{align*}

\item[b)]
From plate scale calculated in a), size of a star for typical seeing 1\,arcsec
\begin{align*}
  &= \frac{1/60~\mathrm{arcmin}}
     {\left(3.44~\mathrm{arcmin/mm} \times 6 \times 10^{-3}~\mathrm{mm}\right)}
     ~\text{pixels}
  && \text{(1 mark)} \\
  &= 0.8~\text{pixel}
\end{align*}
$\Rightarrow$ light from star is mostly contained within
$N_{\mathrm{pix}} = 1$ pixel. \hfill(1 mark)

\emph{Alternative solution:} If a student gets the correct answer for
1\,arcsec $= 0.8$ pixel but clearly explains that the chosen aperture size is
$N \times$ ``seeing'' or a circular aperture with radius $R_{N\mathrm{pix}}$
(where $R_{N\mathrm{pix}}$ is not larger than $2\times$ seeing), calculates
the total number of pixels as the closest round-up integer, and correctly uses
\[ \mathrm{SNR} = \frac{\mathrm{CR} \times t}
   {\sqrt{N_{\mathrm{pix}}\,\mathrm{RON}^2
    + (N_{\mathrm{pix}} \times \mathrm{DN} \times t)}} \]
then the student will be awarded full marks for the relevant parts (this
also applies to part c).

From the definition of the zero-point magnitude given in the question
\[ m = -2.5\log_{10}(\text{count rate}) + \text{Zero-Point} \]
\hfill(2 marks)
\[ m = 21,~\mathrm{ZP} = 18.5 \]
$\Rightarrow$ source count rate
$\mathrm{CR} = 10^{-(21-18.5)/2.5} = 0.1$ per second \hfill(1 mark)

Signal-to-noise ratio
\begin{gather*}
\mathrm{SNR} = \frac{\mathrm{CR} \times t}
  {\sqrt{\mathrm{RON}^2 + (\mathrm{DN} \times t)}} \\
(\mathrm{CR} \times t)^2 = 25(\mathrm{RON}^2 + (\mathrm{DN} \times t))
\end{gather*}
\hfill(1 mark)

Solve quadratic equation
\[ t = \frac{25\,\mathrm{DN}
   + \sqrt{(25\,\mathrm{DN})^2 + 100\,\mathrm{CR}^2\,\mathrm{RON}^2}}
   {2\,\mathrm{CR}^2} \]
\hfill(1 mark)
\[ t = \frac{\dfrac{25}{60}
   + \sqrt{\left(\dfrac{25}{60}\right)^2 + 100\,(0.1)^2\,10^2}}
   {2 \times 0.1^2}
   = 521~\mathrm{s} = 8.68~\text{minutes} \]
\hfill(1 mark)

\item[c)]
If we include the source Poisson noise,
\[ \mathrm{SNR} = \frac{\mathrm{CR} \times t}
   {\sqrt{\mathrm{RON}^2 + (\mathrm{DN} \times t) + \mathrm{CR} \times t}} \]
\hfill(1 mark)

For the Poisson noise dominated case
\[ \mathrm{SNR} \approx \frac{\mathrm{CR} \times t}
   {\sqrt{\mathrm{CR} \times t}} \]
\hfill(1 mark)

Hence
\[ \mathrm{SNR} \propto t^{\frac{1}{2}} \]
\hfill(1 mark)

Re-calculate exposure time
\begin{gather*}
\mathrm{SNR} = \frac{\mathrm{CR} \times t}
  {\sqrt{\mathrm{RON}^2 + (\mathrm{DN} \times t) + \mathrm{CR} \times t}}, \\
(\mathrm{CR} \times t)^2
  = 25(\mathrm{RON}^2 + (\mathrm{DN} + \mathrm{CR})t)
\end{gather*}
\hfill(1 mark)

Solve quadratic equation
\[ t = \frac{25(\mathrm{DN} + \mathrm{CR})
   + \sqrt{25^2(\mathrm{DN} + \mathrm{CR})^2
   + 100\,\mathrm{CR}^2\,\mathrm{RON}^2}}{2\,\mathrm{CR}^2} \]
\hfill(1 mark)
\[ t = 666~\mathrm{s} = 11.1~\text{minutes} \]
\hfill(1 mark)

\item[d)]
In 1\,hr, one can observe $\dfrac{60}{11.1} = 5.4$ \hfill(1 mark)\\
$\Rightarrow$ 5 pointings (round down to the nearest integer). \hfill(1 mark)

Therefore, we need to cover $100~\mathrm{deg}^2$ $\Rightarrow$ one pointing
needs to cover
\[ \frac{100}{(5 \times 4~\text{telescope})} = 5~\mathrm{deg}^2 \]
\hfill(1 mark)\\
or $2.24 \times 2.24$ deg per pointing per telescope. \hfill(1 mark)

Calculate size of each CCD
\begin{align*}
  &= \frac{2.24~\mathrm{deg} \times 60}
     {(3.4~\mathrm{arcmin/mm})(6 \times 10^{-3}~\mathrm{mm})}~\text{pixels}
  && \text{(1 mark)} \\
  &= 6588.2~\text{pixels}
   \;\Rightarrow\; 6589~\text{pixels (pixel number needs to be integer)}
\end{align*}
Therefore, we need minimum $\mathrm{CCD~size} = 6589 \times 6589$ pixels.
\hfill(1 mark)
\end{description}
